\section{The QR Ontology in One Chain}\label{sec:synthesis}
Registered relational objects undergo admissible changes of organization while complete informational accounting is preserved. A scalar complexity readout is derived from a declared object; it does not replace the object or its receipt.
% Original equation number(s) 51; expression unchanged.
\setcounter{equation}{50}
\begin{equation}
\begin{aligned}
\text{relation}&\longrightarrow\text{admissible continuation}\\
&\longrightarrow\text{local transduction}\\
&\longrightarrow\text{propagation and retained receipt}\\
&\longrightarrow\text{completed adjacent handoff}\\
&\longrightarrow\text{composed registered history}\\
&\longrightarrow\text{oriented history readout}.
\end{aligned}
\end{equation}

Mode and phase structure registered presentations without adding handoffs or redefining temporal winding. The common cover supplies a finer carrier presentation, and the relational kernel supplies explicit colored couplings. Neither replaces the completed-event requirement. The scalar-relay equations track complexity changes only after the functional and sufficient local interface are supplied.

A selected executing reference locates the account within this architecture. Its retained context belongs to the complete state rather than to an unspecified outside ledger. Incidence states how that account meets another presentation; reconciliation requires a comparison law that respects the declared transport. These are complementary roles, not three names for the same object.

\section{The Two Branches of Every Propagating Event}
Propagation changes field address; the retained output preserves complementary distinctions. Recovery concerns their joint state. A downstream scalar readout may use less than the complete output, but it must pass the local sufficiency criterion of Section~\ref{sec:scalarrelay}. Selecting a reference does not make unavailable context accessible. Causal reach, relational complexity and temporal winding remain separate readouts.

\section{The Quantum and Geometric Readings}
The history-conference algebra supplies the internal measure correspondence. The geometric program seeks compatible intervals and dynamics from the same event architecture. Each reduction must state the relational distinctions it retains, its native map and its calibrated physical interface. An incidence-mediated comparison may relate different endpoint objects without identifying them. Its existence is a mathematical question about the common action, not a consequence of calling both objects quantum, geometric or oriented.

\section{A Unified Ontological Program}
Carrier, cycle-incidence, conference, outcome-table, modal-frame, census-family, rank-eight kernel and bounded WXYZTI constructions remain separately checked. The new replays add native signed-face incidence transport, the correct realization-kernel naturality, a split signed-axis extension and literal history-phase transport. Imported designations and unexecuted producer claims remain identifiable in the source ledger.

The revised architecture has three distinct construction tasks. The complexity program must select a native functional and a sufficient relay interface. The dynamical program must select an admitted generator, invariant form and represented observables. The orientation program must construct the common incidences and compatible lifted actions. None can be completed solely by selecting a reference convention for the other two.

\section{Conclusion}
Quantum Relativity is a theory of relational complexity governed by conservation of information. The Thalean transducer supplies its local mechanism. Admissible transduction reorganizes relational structure without abandoning the distinctions required for complete accounting. Scalar structure measures a declared aspect of relational organization; phase and Mode describe its organization and orientation; history records composition; and the causal datum bounds how a retained distinction acquires a new address.

An execution nevertheless starts somewhere in its specified apparatus. The executing reference gives the account a situated presentation and retained context without making that position privileged. ``I am 1'' is used as a reference normalization, not as a new physical constant or an absolute center. Another account can select its own reference. The question is then how their registered relations meet and how their descriptions transform together.

Signed incidence provides a concrete answer at an already constructed interface. A traversal sign compares an oriented history with an oriented edge. Shared-incidence products retain relative agreement while cancelling the arbitrary edge-reference signs. More generally, two lifted orientation actions on a common transitive incidence domain admit an equivariant comparison precisely when their isotropy characters agree. The two remaining sign conventions then describe the same covariance structure; selecting either is not a substitute for constructing the domain and actions in the first place.

The finite results enforce that distinction. The 120 bare incidences, the 120-vertex common cover, the $S_5$ sector group and the $A_5\times C_2$ signed operator group cannot be identified by their counts. The history-phase automorphism transports cosets and actors exactly, but does not by itself establish the required state readout or observable coupling. Native chronology, common-adjoint face binding and instrument realization therefore remain open at their actual interfaces.

The retained conservative calculations show how changing organization can preserve a supplied invariant, with rotor and Maxwell examples on separately specified domains. They do not derive those physical models from reference selection or from the carrier. A universal native complexity law and controlled physical limits likewise remain research targets.

The resulting principle is operational: no reference is privileged by the declared presentation laws, but an execution uses a selected reference. The structure supplies the admissible relationships; the executor selects a presentation; the theory must show how the resulting accounts agree. Incidence-based reconciliation makes that final requirement precise without mistaking a convention for a coupling law.
